\scanpage{44}

% B01 | head, item 1
1) Soit $G$ groupe fini d'automorphismes d'un anneau $1/2$ local $O$. $G$ est transitif sur les idéaux maximaux $\underline m$, soit $G_d$ le stabilisateur de $\underline m_0$, et $G_i\subset G_d$ l'ens des élt de $G_d$ qui définissent l'automorphisme identique de $O/\underline m$. Sous ces conditions

% B02 | upper, part a
a) Il existe $x\in O^{G_i}$ tel que $T_{G/G_i}\,x=1$. Donc pour \unsure{tout|$\forall$} $G$-$O$-module $M$, on a

% B03 | upper, stable elements formula
\[\widehat H(G,M)=\widehat H(G_i,M)^{G/G_i}\]

% B04 | next three lines
Donc pour qu'on ait $1\in T_G O$, i.e. $\widehat H(G,M)=0$ ql soit $M$, il f. et s. que $[G_i:e]$ ne soit pas multiple de la caractéristique $p$ de $O/\underline m$.

% B05 | middle, part b
b) Si $G_i=\{e\}$, on a $H^1(G,O^*)=\{e\}$.

% B06 | middle, question beside vertical rule
Plus précisément, quel que soit $G_i$, l'hom naturel $H^1(G,O^*)\longrightarrow H^1(G_i,(\underline O\deleted{\unsure{G|G_i}}/m)^*)$ est bijectif (?). Qu'en est-il dans le cas \deleted{\unsure{G}} des $\operatorname{Gl}(n,\underline O)$ au lieu de $O^*$.

% B07 | left of restriction question
\margin{?}

% B08 | low left, diagonal
% slanted
\margin{dans le cas où \unsure{$O$|$A$} est de dim $1$, car alors $\widehat H(G,M)$!}

% B09 | lower third, after blank
Pour $a$, noter

% B10 | lower third, overlines and sum
\[\begin{aligned}
\overline{\widehat H(G,M)}&=\widehat H(G,\overline M)=\widehat H(G,\bigoplus_{s\in G/G_d}s.\overline M_{\underline m_0})=\widehat H(G_d,\overline M_{\underline m_0})\\
&=\overline{H(G_d,M_{\underline m_0})}.
\end{aligned}\]

% B11 | beneath first equality, brace
\inserted{justification par Zariski; il y a donc hyp. noethérienne \unsure{finie??|fausse??|faire??}}

% B12 | above and left of diagram
Considérons alors le diagramme :

% B13 | foot, left and right vertical maps
\diagram{\begin{tikzcd}[ampersand replacement=\&,column sep=small]
\widehat H(G,M) \arrow[r,"\sim"] \arrow[d,"\wr"'] \& \widehat H(G_i,M)^{G/G_i}_{\deleted{\unsure{G|G_i}}} \arrow[r,"\deleted{\unsure{g|\mathrm{res}}}","\alpha"'] \& \widehat H(G_i,M_{\underline m_0})^{G_d/G_i}\\
\widehat H(G_d,M)^{G/G_d} \arrow[rr,"\beta"] \& \& \widehat H(G_d,M_{\underline m_0})\simeq\widehat H(G_d,M)_{O_d-m\cap O_d}\arrow[u,"\wr"]
\end{tikzcd}}

% B14 | pencil, by enclosed right-hand column
\margin{on passe aux \unsure{complétés|quotients}!}

% B15 | bottom right, upright rule
\margin{$\alpha$ et $\beta$ sont injectifs, à images denses}

% archive sheet number: 43

